\documentclass[a4paper,12pt]{article}
\usepackage[utf8]{inputenc}
\usepackage[T1]{fontenc}
\usepackage{lmodern}
\usepackage{hyperref}
\usepackage{enumitem}
\hypersetup{colorlinks=true}
\title{Asocjacja J1939 w CAN Simulator GUI}
\author{Zespół CAN Simulator GUI}
\date{\today}

\begin{document}
\maketitle

\section{Wprowadzenie}
Od Fazy 14 CAN Simulator GUI łączy interaktywne uczenie asocjacyjne z protokołem SAE J1939. Użytkownik może teraz wskazać zdarzenie lub wartość referencyjną, a algorytm analizuje ramki J1939 na poziomie \textbf{PGN} (Parameter Group Number), \textbf{adresu źródłowego} i konkretnego \textbf{bajtu}.

\section{Działanie}
\begin{enumerate}[noitemsep]
    \item W zakładce \texttt{Uczenie asocjacyjne} wybierz tryb \textbf{J1939} (radiobutton).
    \item Rozpocznij uczenie i oznaczaj zdarzenia (checkbox) lub wprowadzaj wartości referencyjne (np. temperatura).
    \item Algorytm analizuje tylko ramki 29-bitowe (J1939), automatycznie wyodrębnia PGN i adres źródłowy.
    \item Wyniki w tabeli pokazują: \textbf{PGN (nazwa)}, \textbf{ID}, \textbf{Source Address}, \textbf{Bajt}, \textbf{Wartość}, \textbf{Pewność}.
    \item Eksport wzorca (JSON v2) zawiera pełne metadane J1939.
\end{enumerate}

\section{Komponenty}
\begin{itemize}[noitemsep]
    \item \texttt{controllers/j1939\_associative\_controller.py} – dziedziczy po \texttt{AssociativeController}, nadpisuje analizę i dodaje pola J1939.
    \item \texttt{parsers\_j1939.py} – dekompozycja 29-bit ID na PGN, priorytet i adres źródłowy.
    \item \texttt{j1939\_pgn\_definitions.json} – wbudowana baza nazw PGN (ponad 70 wpisów).
\end{itemize}

\section{Testy}
Testy integracyjne znajdują się w \texttt{tests/test\_j1939\_associative.py}. Uruchomienie:
\begin{verbatim}
python3 -m pytest tests/test_j1939_associative.py -v
\end{verbatim}

\section{Ograniczenia}
\begin{itemize}[noitemsep]
    \item Analizowane są wyłącznie ramki z rozszerzonym identyfikatorem (29-bit).
    \item Protokół transportowy J1939 (wieloramkowy) nie jest jeszcze obsługiwany.
\end{itemize}
\end{document}
